\documentclass[11pt]{article}
\usepackage[T1]{fontenc}
\usepackage{amsmath,amssymb,amsthm}
\usepackage[margin=1in]{geometry}

\newtheorem{theorem}{Theorem}
\newtheorem{lemma}{Lemma}
\newtheorem{remark}{Remark}

\title{Laboratory V77: FPT Avoidance from Support Connectivity Branchwidth}
\author{NC0\_k-Avoid Laboratory}
\date{August 2026}

\begin{document}
\maketitle

Let $C:\{0,1\}^n\to\{0,1\}^m$ be an explicitly represented Boolean
circuit whose output gates have fan-in at most three. For a set $S$ of output
gates, define
\[
 \lambda_C(S)=\left|\bigcup_{i\in S}\operatorname{supp}(i)
 \cap \bigcup_{i\notin S}\operatorname{supp}(i)\right|.
\]

\begin{lemma}[Support connectivity]
The function $\lambda_C$ is normalized, symmetric, integer valued,
nonnegative, and submodular. Hence it is a connectivity function.
Furthermore, one oracle evaluation takes $\gamma=O(m)$ time for explicit
fan-in-three supports.
\end{lemma}

\begin{proof}
For each input variable $x$, let $I_x(S)$ be one when $x$ occurs in at least
one gate of $S$ and at least one gate outside $S$, and zero otherwise. The
function $I_x$ is the cut indicator of the set of output gates containing
$x$; it is normalized, symmetric, and submodular. Since
$\lambda_C=\sum_x I_x$, the stated properties follow. Scanning all supports
uses at most three incidences per gate.
\end{proof}

\paragraph{Prior-art decomposition theorem.}
Korhonen and Oum prove that, for a connectivity function on a ground set of
size $m$ given by an oracle of cost $\gamma$, a branch decomposition of width
at most $k$ can be found, or nonexistence certified, in
\[
 2^{O(k)}\gamma m^6\log m+2^{O(k^2)}\gamma m.
\]
The valid simplified upper bound
$2^{O(k^2)}\gamma m^6\log m$ is sometimes used for display. The theorem is
used as a black box and is not reimplemented here.

\begin{theorem}[Support-branchwidth FPT avoidance]
Assume $m>n$, and let $k=\operatorname{bw}(\lambda_C)$. There is an algorithm
that constructs an output word outside $\operatorname{range}(C)$ in time
\[
 2^{O(k)}\gamma m^6\log m
 +2^{O(k^2)}\gamma m
 +O\!\left(m\log m\,A(2k)^2\operatorname{poly}(n,m)\right).
\]
In the explicit support representation, $\gamma=O(m)$. Consequently
$\mathrm{NC}^0_3$-Avoid is fixed-parameter tractable when parameterized by
support connectivity branchwidth.
\end{theorem}

\begin{proof}
Apply the Korhonen--Oum algorithm to $\lambda_C$ to obtain a width-$k$ gate
branch decomposition. Apply the V77 restricted topology-tree transfer to
obtain a rooted binary gate tree of width at most $2k$, height $O(\log m)$,
and external path length $O(m\log m)$. Compile the V75 symbolic residual
circuit on this tree. The V74 exact prefix identity chooses a zero-preimage
child at each output coordinate. Since $m>n$, the final output word has no
preimage. Summing the discovery and avoidance costs gives the bound.
\end{proof}

\begin{remark}
This theorem removes the supplied-decomposition hypothesis from the
parameterized chain. It does not bound $k$ for unrestricted circuits, does not
implement the Korhonen--Oum algorithm, does not prove a standard-model lower
bound, and does not resolve P versus NP. Novelty and peer review are not
claimed.
\end{remark}

\paragraph{References.}
T.~Korhonen and S.-i.~Oum, \emph{Branch-width of connectivity functions is
fixed-parameter tractable}, arXiv:2601.04756, 2026. F.~V.~Fomin and
T.~Korhonen, \emph{Fast FPT-Approximation of Branchwidth}, STOC 2022.

\end{document}
